\section{Exercise 5: Multiplication}
\subsection{Task and method}
Implement \texttt{[yn, yorigin] = multi(x1n, x1origin, x2n, x2origin)}
to calculate the pointwise product
\begin{equation}
  y(n)=x_1(n)x_2(n).
\end{equation}
As in Exercise 4, the input samples must first be aligned by time.
Let $N_1,N_2$ be the input vector lengths and $o_1,o_2$ their origins.
The implementation retains the full common interval
\begin{align}
  n_{\min}&=\min(1-o_1,1-o_2),\\
  n_{\max}&=\max(N_1-o_1,N_2-o_2).
\end{align}
Two zero-filled vectors are created on this interval. Each stored input
sample at time $n$ is inserted at position $n-n_{\min}+1$.
The output is then obtained with the Scilab elementwise operator:
\begin{lstlisting}[numbers=none]
yn = x1aligned .* x2aligned;
yorigin = 1 - nfirst;
\end{lstlisting}
The operator \texttt{.*} multiplies corresponding samples;
\texttt{*} performs matrix multiplication. The pointwise product is
also distinct from the convolution operation in Exercise 6.

\subsection{Source code}
\lstinputlisting[caption={Scilab pointwise multiplication function and example call},label={lst:multi}]{../5.sci}

\subsection{Example and expected result}
Reuse the two input signals from Exercise 4:
\begin{lstlisting}[numbers=none]
[yn, yorigin] = multi([0, 1, 3, -2], 1, [1, 1, 2, 3], 2);
\end{lstlisting}
Both inputs are represented on $n=-1,0,1,2,3$ before multiplication.
\begin{table}[H]
  \centering
  \caption{Time alignment and pointwise multiplication}
  \label{tab:multi}
  \begin{tabular}{lrrrrr}
    \toprule
    Time $n$ & $-1$ & 0 & 1 & 2 & 3\\
    \midrule
    Aligned $x_1(n)$ & 0 & 0 & 1 & 3 & $-2$\\
    Aligned $x_2(n)$ & 1 & 1 & 2 & 3 & 0\\
    Product $y(n)$ & 0 & 0 & 2 & 9 & 0\\
    \bottomrule
  \end{tabular}
\end{table}
The calculated output is
\[
  \texttt{yn}=[0,0,2,9,0],\qquad \texttt{yorigin}=2.
\]
Zero-valued samples at the interval boundaries are retained. The origin
therefore refers to the second position, even though the first nonzero
output occurs at the third position.

\subsection{Execution evidence}
\resultfigure{images/exercise5.png}
  {Aligned input signals and their pointwise product for Exercise 5}
  {fig:multi}
The captured Scilab plot agrees with \Cref{tab:multi}: the output has
two nonzero samples, at $n=1$ and $n=2$.

\subsection{Discussion}
At $n=1$, the product is $1\times2=2$, and at $n=2$ it is
$3\times3=9$. At $n=-1$, $n=0$, and $n=3$, at least one factor is
zero, so the output is zero. The captured output plot therefore contains
two nonzero bars, of amplitudes $2$ and $9$ at $n=1$ and $n=2$,
respectively. In general, a pointwise product can be nonzero only where
both input signals are nonzero.
